
\documentclass[12pt]{article}
\usepackage{amsmath, amssymb}
\usepackage{geometry}
\geometry{margin=1in}
\title{SAT.O3 – Gauge Symmetry from Topological Classes}
\author{SATO.4D Coherence Project}
\date{June 2025}

\begin{document}
\maketitle

\section*{1. Objective}

To derive gauge symmetry groups from the allowed topological classes of filament bundles in 4D, without assuming gauge fields or symmetry a priori.

\section*{2. Allowed Binding Classes}

Filament bundles bind into composite structures with stable topologies:
\begin{itemize}
    \item \( n = 1 \): Unbound filament (e.g., neutrino analog)
    \item \( n = 2 \): Linked pair (meson class)
    \item \( n = 3 \): Borromean or triple link (baryon class)
\end{itemize}

No topologically stable binding exists for \( n \geq 4 \) (per SAT.O4).

\section*{3. Phase-Defined Binding Conditions}

Filament i binds to filament j when:
\[
\phi_i - \phi_j = \frac{2\pi k}{n}, \quad k \in \mathbb{Z}
\]

Phase vector \( \vec{\phi} = (\phi_x, \phi_y, \phi_z) \) defines internal symmetry degrees of freedom.

\section*{4. Emergent Symmetry Group Algebra}

The following gauge groups emerge from phase symmetry:

\begin{itemize}
    \item \( U(1) \): One-filament phase rotations
    \item \( SU(2) \): Two-filament exchange symmetry (spinor doublets)
    \item \( SU(3) \): Three-filament phase permutations (triplet states)
\end{itemize}

These arise from the invariance of binding conditions under respective phase transformations.

\section*{5. Generator Structure}

Each symmetry group has associated generators:
\[
T^a = \text{infinitesimal phase deformation of bundle class}
\]

Commutators reflect underlying filament topological algebra, not Lie algebra imposed by hand.

\section*{6. Topological Origin of Gauge Invariance}

Gauge invariance is not a symmetry of fields, but a redundancy of topological equivalence class:
\[
\gamma_i(\lambda) \sim \gamma'_i(\lambda) \quad \text{iff} \quad \text{same linking class}
\]

Local changes in phase or bundle representation do not alter observable linking.

\section*{7. Module Linkages}

\begin{itemize}
    \item \textbf{O1}: Phase structure defined at filament level
    \item \textbf{O5}: Coupling constants emerge from linking class densities
    \item \textbf{O6}: Gauge symmetry groups define action terms
    \item \textbf{O4}: Restricts allowable symmetry classes to \( n \leq 3 \)
\end{itemize}

\section*{8. Falsifiability Conditions}

\begin{itemize}
    \item Any stable topological structure with \( n > 3 \) falsifies gauge group derivation.
    \item Nonstandard phase-binding patterns must correspond to new symmetries and observable particles.
\end{itemize}

\section*{9. Frame Declaration}

This derivation is conducted in \textbf{Interpretive Mode 1 (True Block)}. All symmetries emerge from phase-aligned binding within topological bundles; no group is assumed.

\end{document}
